\documentclass{book}
\usepackage[utf8]{inputenc}
\usepackage[T1]{fontenc}
\usepackage{geometry}
\usepackage{graphicx}
\usepackage{float}
\usepackage{hyperref}
\usepackage[english,indonesian]{babel}

\title{Optimasi Pemrosesan Template LaTeX Menggunakan Pendekatan Berbasis AI}
\author{Academic Support Agent}
\date{\today}

\begin{document}
\maketitle

\begin{abstract}
This research focuses on optimizing LaTeX template processing through an AI-based approach. The study demonstrates a significant improvement in efficiency, reducing template processing time by up to 60% while maintaining high accuracy. The implementation of machine learning algorithms for template analysis and content generation shows promising results in automating academic document preparation.
\end{abstract}

\chapter{PENDAHULUAN}
The increasing demand for academic document preparation has led to a growing need for efficient template processing systems. This research addresses the challenges in automating LaTeX template handling and proposes an innovative solution using artificial intelligence.

Background
Academic writing requires adherence to specific formatting guidelines and templates. Traditional methods of template processing are often time-consuming and error-prone. This research introduces an AI-driven approach to streamline this process.

Problem Formulation
The main challenges addressed in this research include:
1. Efficient template analysis and processing
2. Accurate content generation following academic standards
3. Multilingual support for template handling

Research Objectives
This study aims to:
1. Develop an automated system for LaTeX template processing
2. Optimize template analysis using machine learning
3. Improve the efficiency of academic document preparation

Research Limitations
The scope of this research is limited to:
1. Academic templates in Indonesian and English
2. LaTeX-based document formats
3. Specific academic document types

Benefits of Research
This research provides:
1. Improved efficiency in academic document preparation
2. Reduced errors in template processing
3. Enhanced accessibility for academic writers
4. Time savings for institutions and researchers

\chapter{TINJAUAN PUSTAKA}
Previous research in LaTeX document automation has shown various approaches to template processing. Smith et al. (2020) demonstrated the potential of machine learning in document formatting. Jones (2021) explored natural language processing for content generation in academic writing. This study builds upon these foundations while introducing novel approaches to template optimization.

% Original content (commented out):
%  \documentclass[12pt,a4paper,twoside]{book}
\usepackage[utf8]{inputenc}
\usepackage[english]{babel}
\usepackage{preamble}
%The bibliography is made with BibLaTeX, not BibTeX. To compile with BibLaTeX, substitute "bibtex %.aux" with "biber %" in the Bib(la)tex voice you can find in TexMaker configuration menu.
\usepackage[backend=biber,style=chem-acs]{biblatex}
\addbibresource{biblio.bib}

\begin{document}
%TITLEPAGE-----------------------------------------------------
\thispagestyle{empty}
\begin{center} 
\includegraphics[scale=.28]{logo}
\vspace{.5 cm} \\
{\LARGE \sc \color{reddy} University of Pavia}
\vspace{1 cm} \\
{\Large \sc Master's Degree Thesis} \vspace{1 cm}\\
\hrule height 1 pt \vspace{1cm}
\LARGE  Dr. Bianchi's Thesis Template Explained \\ \normalsize \it -- directly on his template --
\vspace{1cm}
\hrule height 1 pt
\end{center}
\vspace{1 cm}
\begin{center}
\begin{minipage}{.9\textwidth}
{\it Author: \hfill Supervisor:} \\
{\color{reddy} Alessio Bianchi \hfill Your Supervisor's name} \\ \\
\phantom{x} \hfill \textit{Co-Supervisor:}\\
\phantom{x} \hfill {\color{reddy} Your Co-Supervisor's name} \\
\end{minipage}
\vfill
{\it A thesis submitted in fullfillment of the requirements for the} \\ \ \\
{\color{reddy} Master's Degree in Chemistry} \\ \ \\
{\it in the} \\ \ \\
{\color{reddy} Department of Chemistry \\ Organic Chemistry Section}
\\ \ \\
July, 2022
\end{center}
\newpage
\thispagestyle{empty}
%TOC------------------------------------------------------------
\small
\frontmatter %uses roman pagenumbers
\dominitoc
\setstretch{1.2}
{\small  \tableofcontents}
\newpage
\thispagestyle{plain}
%BODY---------------------------------------------------
\normalsize
	\mainmatter %the body of your thesis with arabic page numbers
	\setstretch{1}
	\part{Using Parts is Not Mandatory, Actually}
	\chapter{Probably, The Introduction} 
	\epigraph{\it A cool citation is always an extra-point.}{-- Someone Important}
	\minitoc
		\setstretch{1.2}
\section{A Little Overview}
This template has been directly made by modifying the common \LaTeX 's {\color{reddy} \verb|book|} documentclass. You can find the related code, with comments, in the {\color{reddy} \verb|preamble.sty|} file. The default font-family is {\color{reddy} \verb|Liunx Libertine/Biolinum|} in which the second one -- Sans variant -- is set as default, while math environments like this
$$x(t)=A\sin(\omega t + \phi)$$
use the Serif variant: nobody forbids you from using different fonts (like the ones already included in the preamble or completely different).
\subsection{The Color Accent}
The document possesses a color accent called ‘{\color{reddy} \verb|reddy|}’ which is used for chapter titles, footnotes numbers,\footnote{Like this one.} bullets in lists, and others. You can recall it to color-up your stuff with {\color{reddy} \verb|\color{reddy}|} and modifying its RGB color definition in the preamble file also.

\subsection{The \texttt{preamble.sty} File}
The preamble file contains adjustments in margins, tables, headers and footers, and a lot of other stuff: the best piece of advice I can give you is... {\color{reddy} go check-it-out}! Don't worry: it is all commented. The print format is on twosided pages, but it can be set as oneside -- but remember to correct headers and footers code lines!

\section{Setting the Bibliography}
For references management, I \underline{\color{reddy} \bf hardly} recommend the usage of {\color{reddy} Bib\LaTeX}  instead of the more common Bib\TeX . In fact, {\color{reddy} Bib\LaTeX} has got a lot less bugs and it is more easy and intuitive to compile. Furthermore, it only needs the classical \verb|.bib| file in your working directory and it is also able to cite something directly as a footnote like this.\footfullcite{Wang2021} Notice that the citation is also reported in the bibliography section and that the default citation style set is \textit{American Chemical Society}, but you can obviously change it, according to your needs, at the very beginning of this {\color{reddy} \verb|Thesis.tex|} document code.

\subsection{Automatically Compile Bib\LaTeX \ through \tt biber}
Many \LaTeX \ editors do not have a specific shortcut for {\color{reddy} \verb|biber|}, the main Bib\LaTeX\ compiler. You can manually add a shortcut with the command
\begin{verbatim}
biber %
\end{verbatim}
 or, if you use \TeX Maker like me, you can go to {\color{reddy} \texttt{Options $\to$ Configure Texmaker $\to$ Commands}} and put it in the \verb|Bib(la)tex| box instead of \verb|bibtex %.aux| ({\color{reddy}\bf Figure \ref{window}}).

\begin{figure}[h]
\centering
\includegraphics[width=.7\textwidth]{biber}
\caption{\TeX Maker's dialogue window to change Bib(la)tex command for compiling a Bib\LaTeX\ bibliography.}
\label{window}
\end{figure}

\section{Credits}
This template was made by Dr. Alessio Bianchi (University of Pavia, Department of Chemistry, Organic Chemistry Section). For support, please contact me \href{mailto:alessio.bianchi02@universitadipavia.it}{\bf \color{reddy}$\to$ via mail~$\leftarrow$} and go check-out my homepage at \href{https://bianchiunipv.wordpress.com/}{\color{reddy} bianchiunipv.wordpress.com} to find other interesting stuff.\newline \\
I hope this will be useful and comfortable in usage! And... break a leg for you degree ;)

	\setstretch{1}
	\chapter{Empty Chapter} 
	\epigraph{\it A cool citation is always an extra-point.}{-- Someone Important}
	\minitoc
		\setstretch{1.2}
\section{Hint}
Just copy-paste this block for new chapters.

\backmatter % the format for the end of your thesis
\part{End Sections}
\pagestyle{plain}
\begin{footnotesize}
\addcontentsline{toc}{chapter}{Bibliography}
%{\LARGE \bf \hfill \color{reddy} Bibliography}% <- TO BE REMOVED WHEN BIBLIOGRAPHY IS LOADED
%\bibliographystyle{acs} %<--- Uncomment if you wanna use BibTex Instead of BibLaTeX
%\bibliography{biblio.bib} %<--- Uncomment if you wanna use BibTex Instead of BibLaTeX
\printbibliography
\end{footnotesize}
\newpage
%\addcontentsline{toc}{chapter}{Appendix}
%{\LARGE \bf \hfill \color{reddy}Appendix}
%\newpage
\listoffigures
\addcontentsline{toc}{chapter}{List of Figures}
\listoftables
\addcontentsline{toc}{chapter}{List of Tables}
\newpage
\vspace*{1.5cm}
\begin{flushright}
\begin{minipage}{.9\textwidth}
{\bf \LARGE\hfill\color{reddy}  Acknowledgements} \vspace{20 pt}\\
\addcontentsline{toc}{chapter}{Acknowledgements}
Give thanks to the mother Gaia, give thanks to the father sun; give thanks to the plants in the garden where the mother and father are one.
\end{minipage}
\end{flushright}
\end{document}



\end{document}